\section{Two improved twelve-positive-variable constructions}
\label{sec:positive-improved}

There are two different improvements.  The economical construction evaluates
exactly package 30's degree-$1081$ polynomial in its four coefficient ports
in $108$, instead of $111$,
binary arithmetic operations.  Its three-equation system costs $104$
operations, or $103$ after a justified shift of the positive transport
coordinate.  The second construction changes the polynomial: it has degree
$1065$ in the same twelve positive unknowns, compared with the previous
degree record $1069$.  It costs $118$ operations; its three-equation versions
cost $114$ and $113$ operations, with maximum residual degree $1056$.
The latter combines package 24's sharp weight, package 25's exclusion of
the negative main-root branch, and package 30's quadratic-remainder norm.
The smaller carrier permitted by the branch exclusion creates two additional
leading cancellations.  These are comparisons with the frozen repository
collection, not assertions of a global minimum.

All results in this section use one complete coefficient tuple satisfying
the native blank-compiler interface.  The fixed ports are
$\beta=B-1,K_0,M_C,M_F$, with $M_F$ the shifted high mask.  They are fixed
integers, not additional unknowns.  The coefficient construction elsewhere
in this report supplies the tuple.  Changing only one port is not justified.

\subsection{Common polynomial definitions}

The supplied variables are exactly
\[
 J,F,\alpha,z,h,i,T,s,w,\eta,\zeta,t\in\mathbb Z_{>0}.
\]
Every other displayed symbol denotes a polynomial expression in these
variables.  Put
\begin{align*}
 q&=\beta J+1,& C_w&=q-F-\alpha,&X&=wq,&Y&=sq^3,&E&=XY,\\
 k&=\eta+\zeta,&c&=kY+\eta,&a&=E+Y,&H&=4a+3,&\Delta&=a^2+H,\\
 \ell&=EkY,&A&=1+\ell(\ell+k),&G&=1+\Delta c^2,&b_0&=X+ac.
\end{align*}
In particular $\Delta=(a+2)^2-1$.  The unchanged index and transport
expressions are
\begin{align*}
 R&=((q-1)(q-F)+\alpha)(q^2-1)+(M_C+qM_F)J,\\
 N_k&=k-hE-R,\\
 N_t&=(K_0+w)C_w+q-F-z(q-1).
\end{align*}
Define the strong and auxiliary expressions
\begin{align*}
 I&=ic^2,&D&=1+\Delta I^2,&Q&=\Delta^2 I^2,&d&=Q-1,\\
 v&=cT,&b&=c+RD,&e&=Q(Dv^2+b^2)-1,&B_a&=2Qvb.
\end{align*}
Notice that $Q=(\Delta I)^2$ is the square needed for rational Pell
saturation; it is not an arbitrary positive auxiliary parameter.

Select one of two pairs $(M,W)$, in both cases putting $L=1+HM$:
\begin{align}
 M_{\rm econ}&=k(I^2+T^2),&
 W_{\rm econ}&=b_0+Ht, \label{eq:economical-new}\\
 C_*&=E+2(X+Y)+sq^2, \nonumber\\
 M_{\rm canc}&=I\bigl(2T+RC_*ic\bigr),&
 W_{\rm canc}&=Lb_0+Ht. \label{eq:cancellation-new}
\end{align}
For a formal variable $x$ let
\begin{align*}
 C(x)&=d(x^2+M^2A)-e,\\
 F_*(x)&=\bigl[C(x)^2+DB_a^2-4d^2x^2M^2A\bigr]^2
                    -4DB_a^2C(x)^2,\\
 F_H(u)&=H^8F_*(u/H),\\
 N&=F_H(W-L\sqrt G)F_H(W+L\sqrt G).
\end{align*}
The last expression is an integer polynomial.  Its radicals describe a
finite conjugate product and are not part of the arithmetic language or
the list of quantified unknowns.  A complete division-free evaluator is
given below.  The product equation and unshifted system are
\begin{equation}
 P=N_kN_t(1+N)-1=0,
 \qquad\text{and}\qquad N_k=1,\quad N_t=1,\quad N=0.
 \label{eq:two-positive-presentations}
\end{equation}

\subsection{A cheaper exact norm evaluator}

The following scalar formulas define $N$ without radicals, division, or
resultants.  They are also the basis of the literal JSON gate lists:
\begin{align*}
 m&=HM,&n&=m^2A,&d_n&=dn,&Z&=H^2Q,\\
 P_H&=ZDv^2,&U_H&=Zb^2,&e_H&=P_H+U_H-H^2,&j&=4P_HU_H,\\
 x_0&=d(W^2+L^2G),&\mathcal T&=2x_0,&
 \kappa&=\bigl[d(W^2-L^2G)\bigr]^2,\\
 \mathcal A&=2(d_n+e_H),&
 \mathcal B&=(d_n-e_H)^2-j,\\
 a_r&=\mathcal T-\mathcal A,&b_r&=\mathcal B-\kappa,&
 c_r&=16d_nj,\\
 v_r&=a_r(a_r\mathcal T+2b_r)-c_r,&
 u_r&=b_r^2-a_r^2\kappa,\\
 N&=u_r(u_r+\mathcal T v_r)+\kappa v_r^2.
\end{align*}
Here $d_n=dn$ means the product of $d$ and $n$; it introduces no new
parameter.  In the emitted circuit $v^2=c^2T^2$, so the separately useful
$T^2$ is computed once.

To prove the evaluator, the elementary identity
\[
 (C^2+j-K)^2-4jC^2=(C^2-K-j)^2-4Kj
\]
rewrites $F_H(u)$, with $x=du^2$, as
\[
 (x^2-\mathcal A x+\mathcal B)^2-c_rx.
\]
For $u=W\pm L\sqrt G$, the two corresponding $x$ values have sum
$\mathcal T$ and product $\kappa$.  Reduction modulo
$x^2-\mathcal T x+\kappa$ first replaces
$x^2-\mathcal A x+\mathcal B$ by $a_rx+b_r$, and then replaces its
square minus $c_rx$ by $u_r+v_rx$.  Its norm is exactly
$u_r^2+\mathcal T u_rv_r+\kappa v_r^2$.  This argument works when the
radicands are squares or zero and when conjugates coincide; no
discriminant is divided out.

Three reassociations each save one multiplication relative to package 30:
compute the common factor $H^2Q$ before forming $P_H,U_H$; compute
$\kappa$ by multiplying before squaring; and use the displayed Horner
form for the final norm.  These are identities over every commutative
ring.  At the same coefficient tuple, the economical polynomial is therefore
unchanged, including all its integer zeros.  A different complete tuple
produces a different fixed polynomial through the same compiler interface.
The sharper construction changes $M,L,W$ as stated;
its new degree is proved below.

\subsection{Weak signs and the inherited interface}

We use three precisely identified consequences of the blank-compiler
theorem.  First, the weak index/transport bootstrap, valid for
$N_k=N_t=\epsilon\in\{1,-1\}$, supplies
\begin{equation}
 \begin{gathered}
 C_w\ge0,\quad q\ge16,\quad X\ge q,\quad Y\ge q^3,\\
 3q<R<q^4-2,\quad E\ge q^4>R+2,\quad
 k=R+\epsilon+hE>E,\quad kY<c<k(Y+1).
 \end{gathered}\label{eq:positive-weak-bootstrap}
\end{equation}
Second, restoring the normalized first, main, strong and auxiliary Pell
equations, together with the main congruence, permits the inherited signed
auxiliary-quotient soundness theorem.  It forces $\epsilon=1$ and extracts
a blank K215 halt.  Third, every blank halt has a canonical positive
normalized witness.  Those canonical witnesses satisfy $C_w>0$ and
$X=2^R$.  These are semantic compiler dependencies, not consequences of
checking four integers modulo a few small primes.

The hypotheses of the weak bootstrap follow from a root of the product
equation before any radical is recovered.  Indeed
\[
 \kappa=x_0^2-4d^2W^2L^2G,
 \qquad N\equiv(u_r+x_0v_r)^2\pmod4.
\]
At a zero of $P$, the three integer factors $N_k,N_t,1+N$ are units.
The last cannot equal $-1$, since this would give $N=-2$, which is not a
square modulo $4$.  Thus $N=0$ and $N_k=N_t=\epsilon$.  For a solution
of the system, $\epsilon=1$ directly.  No consistency assumption is used.

\subsection{All branches and the sharp separation bound}

Factoring the finite norm shows that $N=0$ supplies coherent roots
$g,\tau,f,y$, initially in a finite algebraic extension, with
\begin{equation}
 \begin{gathered}
 g^2=G,\quad \tau^2=A,\quad f^2=D,\\
 dy^2=e-B_af=Q(vf-b)^2-1,\\
 W=Lg+H(M\tau+y).
 \end{gathered}\label{eq:coherent-positive-roots}
\end{equation}
Conversely any such coherent choice makes a factor of $N$ zero.  This
factorization neither divides by a radical nor presupposes the rationality
of any root.  The bounds that follow apply to every complex conjugate of
$y$; such a conjugate can be imaginary.

For the economical weight, put $A_0=a+2$.  From
\eqref{eq:positive-weak-bootstrap},
$A_0<2E$, $c<2kY$, $D<A_0^2I^2<4E^2I^2$, and $c<I^2$.
Every auxiliary branch satisfies
\[
 |y|<2(v\sqrt D+b).
\]
The two elementary estimates
\[
 2v\sqrt D<8kEYIT<\tfrac12M_{\rm econ}kY^2,
 \qquad
 2b<10RE^2I^2<\tfrac12M_{\rm econ}kY^2
\]
use $I^2+T^2\ge2IT$, $kY>8E$, $k>E$, and $Y^2>20R$.
Moreover
\[
 \sqrt A>\ell=EkY>kY^2,
 \qquad \sqrt G-\sqrt A>k(Y^2+Y-1)>kY^2.
\]
For the second inequality use
$\sqrt G>(a+1)c>(E+Y+1)kY$ and
$\sqrt A<\ell+k$.  We obtain the two uniform bounds
\begin{equation}
 |y|<M_{\rm econ}\sqrt A,
 \qquad |M_{\rm econ}\tau+y|<M_{\rm econ}\sqrt G.
 \label{eq:economical-separation}
\end{equation}

For the cancellation weight a tighter, differently shaped estimate is
needed.  Let $S=\Delta I$, so $Q=S^2$.  The weak bootstrap implies
$\Delta>R+2$ and $I^2\ge c^4$, hence
\[
 Q-b=\Delta I^2(\Delta-R)-c-R>0.
\]
Consequently $b\le d$ and $\sqrt D<S<d$.  For
$Z_1=v\sqrt D+b>1$, every auxiliary branch satisfies
\[
 |y|^2\le\frac{QZ_1^2+1}{d}
 <\left[Z_1\left(1+\frac1d\right)\right]^2.
\]
The strict inequality follows on subtracting the left member from the
right: the numerator is $QZ_1^2-d>0$.  Therefore
\begin{equation}
 |y|<Z_1+v+1<cT((a+2)I+1)+c+RD+1.
 \label{eq:sharp-y-new}
\end{equation}
Set $\rho=(1+1/Y)/E$.  Because $c/\ell<\rho$ and
$(a+3)\rho<2$, the $T$ term in \eqref{eq:sharp-y-new}, divided by
$\sqrt A>\ell$, is less than $2TI$.  The remaining large term is
controlled by the exact rational identity
\[
 (a+2)^2\rho
 =E+2Y+X+6+\frac YX+\frac5X+\frac4Y+\frac8E+\frac4{EY}.
\]
Here $Y/X=sq^2/w\le sq^2$ and
$5/X+4/Y+8/E+4/(EY)<1$.  Since $X\ge16$, this gives
$(a+2)^2\rho<C_*-2$.  Also $(c+1)/\ell<2$.  Thus
\[
 \frac{|y|}{\sqrt A}
 <2TI+R(C_*-2)i^2c^3+2
 \le 2TI+RC_*i^2c^3=M_{\rm canc}.
\]
We have proved, without recovering any square root,
\begin{equation}
 |y|<M_{\rm canc}\sqrt A,
 \qquad |M_{\rm canc}\tau+y|<2M_{\rm canc}\sqrt A.
 \label{eq:cancellation-separation}
\end{equation}

Finally, the same weak bounds imply $G>A$ and $b_0>\sqrt A$.  Indeed
$c>kY$ and $\Delta>E^2+2EY$ give
\[
 \Delta c^2>k^2(E^2Y^2+EY)=A-1.
\]
For the second inequality $b_0>akY$ and
\[
 a^2k^2Y^2-A=k^2\{EY(2Y^2-1)+Y^4\}-1>0.
\]

\subsection{Why the smaller carrier is sound}

In the economical construction, the two arguments of $F_*$ associated
with the two main roots are $(W\pm L\sqrt G)/H$.  If $G$ were nonsquare,
the vanishing of one factor of its rational quadratic norm would force
the other to vanish as well.  But the two arguments have distance
$2L\sqrt G/H>2M\sqrt G$, whereas each root of $F_*$ has modulus less
than $M\sqrt G$ by \eqref{eq:economical-separation}.  This contradiction
proves that $G$ is a square.  The selected $g$ is positive: if it were
negative then \eqref{eq:coherent-positive-roots} would give
$W<-L\sqrt G+HM\sqrt G<0$, contrary to $W=b_0+Ht>0$.

For the cancellation construction, the weighted carrier makes the same
conclusion possible with only $L=1+HM$, instead of package 24's
$1+2HM$.  The branch $g=-\sqrt G$ would require
\[
 x=\frac{W+L\sqrt G}{H}
   =t+\frac LH(b_0+\sqrt G)>2M\sqrt A,
\]
because $b_0>\sqrt A$, $\sqrt G>\sqrt A$, and $L/H>M$.
This contradicts \eqref{eq:cancellation-separation}.  Thus the negative
branch is impossible.  If $G$ were nonsquare, rational conjugacy would
force that branch to vanish whenever the positive branch did, again a
contradiction.  Hence $g=\sqrt G$ is a positive integer.

In either construction, work in a finite normal splitting field containing
the roots in \eqref{eq:coherent-positive-roots}.  Now $g$ is rational.
If $\tau$ were irrational, an automorphism would send it to $-\tau$ while
fixing $g$.  Subtracting the two linear relations would require a change
of magnitude $2M\sqrt A$ in $y$.  Every conjugate of $y$ has magnitude
strictly less than $M\sqrt A$, an impossibility.  Thus $\tau$ is rational,
and hence an integer because $\tau^2=A\in\mathbb Z$.
The linear relation makes $y$ rational, and
$f=(e-dy^2)/B_a$ is rational because $B_a>0$.
Since $f^2=D\in\mathbb Z$, $f$ is integral.

\paragraph{Rational Pell saturation.}
If integers $S\ge2,V$ and rational $y$ satisfy
$S^2V^2-(S^2-1)y^2=1$, then $y$ is integral.
For completeness write $y=u/v_0$ in lowest terms, $v_0>0$.
Its equation gives $v_0^2\mid S^2-1$.
The positive nonsquare integer $D_0=(S^2-1)/v_0^2$ has the two
integer-coordinate Pell solutions $(S,v_0)$ and $(S|V|,|u|)$.
Both coordinates in each pair are nonzero.  Write these pairs as
$(\chi_j,\psi_j)$ and $(\chi_l,\psi_l)$ in powers of the fundamental
positive integer-coordinate Pell unit.  The addition formulas give
$\chi_{r+2j}\equiv-\chi_r\pmod{\chi_j}$ and
$\chi_{2j-r}\equiv-\chi_r\pmod{\chi_j}$.
Because $\chi_j\mid\chi_l$, reducing $l$ modulo $2j$ and reflecting
into $[0,j]$ shows that $l$ is an odd multiple of $j$; strict growth
excludes every index smaller than $j$.  The Pell multiplication formulas
then give $\psi_j\mid\psi_l$, so $v_0\mid u$ and $v_0=1$.
This uses the integer-coordinate Pell group, not a larger group of
half-integral units.

Apply the lemma with $S=\Delta I$ and $V=vf-b$.  It proves $y\in\mathbb Z$.
Since $L\equiv1\pmod H$ and in both variants $W\equiv b_0\pmod H$,
the integral linear relation now yields $g\equiv b_0\pmod H$.
Therefore $\gamma=(g-b_0)/H$ is integral.  It is also positive:
$g>(a+1)c$ and $c>k>E>X$ imply
$H\gamma=g-ac-X>c-X>0$.

We have restored exactly
\begin{align*}
 \tau^2-\ell(\ell+k)&=1,\\
 (X+ac+H\gamma)^2-\Delta c^2&=1,\\
 f^2-\Delta i^2c^4&=1,\\
 Q[c(Tf-1)-Rf^2]^2-(Q-1)y^2&=1,
 \qquad N_k=N_t=\epsilon.
\end{align*}
The auxiliary identity uses $f^2=D$ and
$vf-b=c(Tf-1)-Rf^2$.  Absolute values may be taken for $\tau$ and $y$.
If $f<0$, change the signs of $f$ and $T$ together, preserving $Tf$.
The inherited signed-quotient compiler theorem applies, forces
$\epsilon=1$, and recovers a blank K215 halting computation.

Conversely, take a canonical positive normalized compiler witness, with
$g=b_0+H\gamma$ and positive $\tau,f,y$.
For the economical construction choose
\[
 t=\gamma+M(g+\tau)+y>0;
\]
for the cancellation construction choose
\[
 t=L\gamma+M\tau+y>0.
\]
The coherent factor then vanishes and $N_k=N_t=1$.
Every canonical witness can be extended in this way; no new restriction
on its strong Pell index is imposed.
This proves both directions of existential equivalence with the authentic
compiler.  Composing with the K215 syntactic-refutation calibration and
negating existential solvability proves, within ZFC, that absence of
positive integer solutions is equivalent to $\operatorname{Con}(\mathrm{ZFC})$.

\subsection{The shifted systems}

In the system only, replace the transport equation by
\begin{equation}
 (K_0+w)C_w=F+z(q-1).
 \label{eq:shifted-positive-system}
\end{equation}
Here the new positive variable is related to the old one by
$z_{\rm old}=z+1$.  Every shifted solution therefore gives an unshifted
solution, so soundness is immediate.  For completeness use a canonical
compiler witness, for which $C_w>0$, $X=wq=2^R$, $R>3q$, and $F<q$.
Consequently $w>q>F$, and the original transport equation gives
\[
 (z_{\rm old}-1)(q-1)=(K_0+w)C_w-F>0.
\]
Thus $z_{\rm old}\ge2$ and $z=z_{\rm old}-1$ is positive.
This positivity is proved after restoring the compiler, rather than
inserted into the weak bootstrap.  Directly comparing the two sides of
\eqref{eq:shifted-positive-system} saves one addition/subtraction.
The product equations retain the unshifted variable.

\subsection{Exact degrees, including the new cancellation}

All four coefficient ports have degree zero and all twelve supplied
variables have degree one.  The economical norm is the unchanged package
30 polynomial, so its exact degree is $1072$; its product degree is $1081$.
Its leading form, writing $q_0=\beta J$ and $k_0=\eta+\zeta$, is
\[
 (N_{\rm econ})_{1072}
 =2^{32}i^{48}k_0^{96}w^{80}s^{176}q_0^{656}(q_0-F)^{16}.
\]

Here is a full upper and lower degree argument for the new construction.
Replace every supplied variable by $\lambda$ times itself, and use a
subscript $0$ for its highest homogeneous part.  In the rational function
field of the unscaled variables, the relevant square roots have formal
Laurent expansions at $\lambda=\infty$.  Their leading radicands are
nonzero squares, so the ordinary formal binomial expansion constructs all
branches without fractional orders.  The basic degrees are
\[
 \begin{array}{c|rrrrrrrrrr}
  \text{expression}&E,H&Y&c&R&I&d&b&A,G&M&L\\ \hline
  \text{degree}&6&4&5&4&11&46&38&22&27&33.
 \end{array}
\]
Define three nonzero leading forms
\[
 \mathfrak B=R_0E_0^2i^2c_0^4,\qquad
 \mathfrak D=R_0E_0Y_0i^2c_0^4,\qquad
 \mathfrak C=R_0E_0(sq_0^2)i^2c_0^4.
\]
Their degrees are $38,36,35$, respectively.

The first radical satisfies $\sqrt A-\ell=O(\lambda)$.
For every $f$ branch the auxiliary roots satisfy
\[
 y=\varepsilon(b-vf)+O(\lambda^{-8}),\qquad
 \varepsilon\in\{1,-1\},\qquad\deg(vf)=23.
\]
The exact identity
\[
 C_*\ell-\Delta c
 =(2EX+Esq^2-Y^2-4E-4Y-3)c-EC_*\eta
\]
has unique highest homogeneous term $E_0(sq_0^2)c_0$, of degree $14$.
Consequently
\begin{equation}
 M\sqrt A-b=\mathfrak C+\text{terms of order below }35.
 \label{eq:new-cancellation-first}
\end{equation}
Indeed the difference is
$Ri^2c^3(C_*\ell-\Delta c)+O(\lambda^{28})$.

Let $g_+=\sqrt G$ and $\gamma_\pm=(\pm g_+-b_0)/H$.
The small quotient root $\gamma_+$ has order $-1$, hence
$L\gamma_+$ has order $32$.
Also $g_+=ac+O(\lambda^5)$ and
\[
 g_++b_0-2\sqrt A
   =2Y_0c_0+\text{terms of order below }9.
\]
Multiplying by $M$, whose top form is $R_0E_0i^2c_0^3$, gives
\begin{equation}
 M(g_++b_0-2\sqrt A)
   =2\mathfrak D+\text{terms of order below }36.
 \label{eq:new-cancellation-second}
\end{equation}

The exact sixteen-branch factorization is
\begin{equation}
 N=d^8H^{16}\prod(t-y-M\tau-L\gamma),
 \label{eq:new-sixteen-factor}
\end{equation}
where the coherent choices of $g,\tau,f,y$ are all included with
multiplicity.  For $\gamma_+$, four factors have degree $35$ with signed
top forms $\pm\mathfrak C$; the other four have degree $38$ with signed
top forms $\pm2\mathfrak B$.
For $\gamma_-$, the two branches with $\tau=+\sqrt A$ and
$y=+(b-vf)+O(\lambda^{-8})$ have degree $36$ and top form
$2\mathfrak D$, by \eqref{eq:new-cancellation-first} and
\eqref{eq:new-cancellation-second}.
The remaining six have degree $38$: four top forms are $2\mathfrak B$
and two are $4\mathfrak B$.
This exhausts every branch.  It proves the degree upper bound, not just
a degree observed on a numerical specialization.  Multiplying the
nonzero top forms gives
\begin{align*}
 \deg N&=8\cdot46+16\cdot6+4\cdot35+2\cdot36+10\cdot38=1056,\\
 N_{1056}
 &=2^{14}d_0^8H_0^{16}\mathfrak B^{10}\mathfrak C^4\mathfrak D^2\\
 &=2^{46}i^{48}k_0^{96}w^{74}s^{176}q_0^{646}(q_0-F)^{16}.
\end{align*}
This form is nonzero for every fixed $\beta>0$.
Since
\[
 (N_k)_7=-hwsq_0^4,\qquad
 (N_t)_2=w(q_0-F-\alpha)-zq_0,
\]
the product has exact degree $1056+7+2=1065$.
For example its coefficient of
\[
 J^{667}hz\,i^{48}\eta^{96}w^{75}s^{177}
\]
is $2^{46}\beta^{667}>0$.
The three-equation residual degrees are exactly $7,2,1056$; shifting
$z$ by one preserves these leading degrees.

\subsection{Literal cost, checks, and proof status}

The supplied evaluator charges every binary $+,-,\times$, including one
by a fixed numeral.  Intermediate expressions are shared; comparisons
in systems are free under the repository convention.
\begin{center}
\begin{tabular}{lrrrr}
\hline
Construction & Equations & Degree & Multiplications & Total operations\\
\hline
Economical product &1&1081&65&108\\
Economical system &3&1072&63&104\\
Economical shifted system &3&1072&63&103\\
Cancellation product &1&1065&72&118\\
Cancellation system &3&1056&70&114\\
Cancellation shifted system &3&1056&70&113\\
\hline
\end{tabular}
\end{center}
The checks compare each new norm with an independently implemented
eight-dimensional multiquadratic conjugate product on $50$ signed integer
assignments, check the parity and transport identities on those assignments,
and reject dead gates or unused ports.  Full univariate coefficient lists
over two primes confirm the cancellations and the stated leading forms;
this is stronger than merely propagating degree upper bounds through the
gates.  The mathematical arguments above, rather than those finite
examples, establish the unbounded root-recovery theorem and uniform exact
degrees.  No new end-to-end Lean or Rocq check is claimed, and no assertion
that ZFC is consistent is made.
